% =============================================================================
%  DESARROLLO
% =============================================================================
\section{Desarrollo}

\subsection{La IA no destruye solo empleos: destruye sentido}

Existe una narrativa tranquilizadora que se repite como mantra cada vez que la tecnología desplaza trabajos: ``la revolución industrial también eliminó empleos y al final creó más de los que destruyó''. Esta analogía, sin embargo, ignora una diferencia fundamental. Las revoluciones anteriores mecanizaron la fuerza física; la revolución de la IA mecaniza la cognición. Por primera vez en la historia, la máquina no solo reemplaza lo que nuestros músculos hacen, sino lo que nuestras mentes producen. Y a diferencia del obrero del siglo~XIX que podía reconvertirse en operador de máquina, el trabajador del siglo~XXI desplazado por la IA se enfrenta a la pregunta más desoladora que puede hacerse un ser humano: ¿para qué sirvo?

Los investigadores \textcite{frey2013} estimaron que el 47\,\% de los empleos en Estados Unidos tenían alta probabilidad de automatización en las siguientes dos décadas. Una década después, con la llegada de los modelos de lenguaje de gran escala, esa cifra se antoja conservadora. No se trata solo de empleos rutinarios: se trata de diagnósticos médicos, redacción periodística, análisis legal, generación de código. Se trata, irónicamente, de los trabajos que se consideraban seguros precisamente porque requerían inteligencia.

Lo que más inquieta no es el número de empleos perdidos, sino lo que esos empleos representaban. Para millones de personas, el trabajo no es solo un medio de subsistencia: es identidad, propósito, estructura del tiempo, fuente de pertenencia social. Cuando una empresa automatiza un área, no desaparece únicamente el ingreso; desaparece la razón por la que alguien se levantaba a las cinco de la mañana, el orgullo del programador en el sistema que construyó, la dignidad del contador que resolvía el problema que nadie más podía. La IA no destruye únicamente empleos: destruye, en muchos casos, el sentido que las personas construyeron alrededor de ellos. Y eso, a diferencia de un salario, no se puede reemplazar con un subsidio.

\subsection{La tecnología tiene política: el ingeniero en sistemas como actor social}

Desde la perspectiva de los estudios de Ciencia, Tecnología y Sociedad (CTS), la tecnología nunca es neutral. El politólogo \textcite{winner1980} planteó que los artefactos tecnológicos ``tienen política'': incorporan en su diseño las relaciones de poder, los valores y los intereses de quienes los crean. La Inteligencia Artificial no es la excepción. Cada sistema de IA refleja las prioridades de quienes lo financian, y quienes financian la gran mayoría de estos sistemas son corporaciones cuyo objetivo explícito es la maximización del valor para el accionista. No es una conspiración; es simplemente la lógica del sistema económico en el que operamos.

La globalización intensifica esta dinámica. Los beneficios económicos de la automatización no se distribuyen entre quienes pierden sus empleos, sino que se concentran en las corporaciones transnacionales ---principalmente radicadas en unos pocos polos tecnológicos del mundo--- que desarrollan y comercializan estos sistemas. El \textcite{wef2020} estimó que para 2025 la automatización destruiría 85 millones de puestos de trabajo a nivel global mientras creaba 97 millones nuevos. La diferencia parecería positiva sobre el papel, pero esos 97 millones de empleos nuevos requieren competencias altamente especializadas, acceso a educación de calidad y conectividad: recursos distribuidos de manera profundamente desigual. El resultado no es un mundo con más empleo: es un mundo donde la brecha entre quienes pueden acceder a la economía digital y quienes no, se vuelve abismal.

Como ingeniero en sistemas, esta realidad me coloca en una posición incómoda que no puedo ignorar. Yo soy parte de la maquinaria que produce esta transformación. Cada línea de código que escribo, cada sistema que diseño, cada solución que optimizo forma parte de un proceso que tiene consecuencias sobre personas reales. Esta incomodidad es precisamente la que el Nuevo Humanismo quiere que sintamos, porque es el punto de partida de una reflexión ética genuina. \textcite{han2012} describe nuestra época como una \textit{sociedad del rendimiento}, donde el imperativo de la eficiencia ha reemplazado al imperativo del sentido. Un ingeniero formado únicamente en rendimiento ---en cuánto procesa, cuánto optimiza, cuánto escala--- es exactamente el tipo de profesional que esa sociedad produce y reproduce. Un ingeniero formado también en humanismo es, quizás, el que puede interrumpir el ciclo.

El problema es que nuestra formación académica raramente nos prepara para esa interrupción. Aprendemos estructuras de datos, algoritmos, paradigmas de programación, arquitecturas de software. Aprendemos a resolver problemas técnicos con elegancia y eficiencia. Lo que rara vez aprendemos es a preguntarnos si el problema que estamos resolviendo vale la pena resolverlo, o a quién beneficia realmente la solución. Como señala \textcite{borgmann1984}, la tecnología moderna actúa como un ``dispositivo'' que oculta su mecanismo y entrega únicamente su producto final, separando al usuario ---y al mismo creador--- de la complejidad y del costo real de lo que produce. La IA es, en muchos sentidos, el dispositivo definitivo: genera valor sin mostrar el costo humano, produce eficiencia sin preguntar a quién le cuesta esa eficiencia.

\subsection{El Nuevo Humanismo como respuesta y como exigencia}

El Nuevo Humanismo, tal como lo concibe la \textcite{unesco2015}, no es una nostalgia por un mundo sin tecnología: es una convicción de que el desarrollo humano debe ser el centro y el fin de toda transformación social, incluida la tecnológica. En ese marco, la pregunta no es si debemos desarrollar IA ---esa decisión ya fue tomada por fuerzas que ningún individuo controla--- sino cómo debemos desarrollarla, bajo qué principios, con qué límites y en función de qué valores.

\textcite{harari2016} plantea un escenario tan probable como perturbador: una humanidad dividida entre una pequeña élite de personas relevantes para la economía algorítmica y una vasta mayoría que simplemente ha quedado obsoleta. No se trata de ciencia ficción distante; se trata de una extrapolación razonable de las tendencias actuales. Y lo que hace que este escenario sea tan profundamente triste no es la pobreza material que implicaría, sino la pérdida de propósito. Una persona sin un rol reconocido en la sociedad, sin la experiencia de ser necesaria, no es solo una persona pobre: es una persona sin sentido.

El Nuevo Humanismo responde a esto con una afirmación radical: el valor del ser humano no puede depender de su productividad económica. Esta afirmación, que parece obvia cuando se enuncia, es subversiva dentro del sistema de valores del mercado global. Nos dice que hay que diseñar tecnología no para maximizar la rentabilidad de las corporaciones, sino para ampliar las posibilidades de la vida humana; para distribuir el ocio liberado por la automatización en lugar de concentrar la riqueza que genera; para emplear la eficiencia de las máquinas en resolver los problemas que el mercado nunca resolverá porque no son rentables: el cambio climático en comunidades pobres, el acceso a la salud en regiones marginadas, la educación de calidad para quienes no pueden pagarla.

Este es el reto que tenemos como ingenieros en sistemas: no somos solo técnicos que ejecutan especificaciones. Somos actores sociales con una capacidad real de influir en cómo se despliega esta tecnología. Y esa capacidad viene acompañada de una responsabilidad que no podemos tercerizar a los gerentes, a los inversores o al mercado.
